\documentclass[11pt,a4paper]{article}
\usepackage[utf8]{inputenc}
\usepackage[T1]{fontenc}
\usepackage{amsmath, amssymb, amsthm}
\usepackage{geometry}
\geometry{margin=1in}

\title{04 - Erd\H{o}s-Gy\'{a}rf\'{a}s Conjecture: A Partial Proof Sketch}
\author{Charles EDOU NZE\thanks{Charles EDOU NZE, chercheur ind\'{e}pendant}}
\date{}

\newtheorem{theorem}{Theorem}
\newtheorem{lemma}{Lemma}
\newtheorem{definition}{Definition}

\begin{document}

\maketitle

\begin{abstract}
This document outlines a structural analysis and a partial proof sketch for the Erd\H{o}s-Gy\'{a}rf\'{a}s conjecture, stating that every graph with minimum degree 3 contains a simple cycle whose length is a power of 2. We present rigorous axiomatic definitions, contextual literature review, lemma decomposition, and a detailed informal proof strategy aimed at subsequent formalization in Lean 4.
\end{abstract}

\section{Definitions and Axiomatic Setup}
\begin{definition}
Let $G = (V, E)$ be a finite, simple, undirected graph. The minimum degree of $G$, denoted $\delta(G)$, is defined as $\min_{v \in V} \deg(v)$.
\end{definition}

\begin{definition}
A cycle $C$ in $G$ is a sequence of distinct vertices $v_1, v_2, \ldots, v_k$ such that $(v_i, v_{i+1}) \in E$ for $1 \leq i < k$, and $(v_k, v_1) \in E$. The length of $C$, denoted $\ell(C)$, is $k$.
\end{definition}

\section{Literature Context}
The Erd\H{o}s-Gy\'{a}rf\'{a}s conjecture, posed in 1995, remains open. It relates to the existence of cycles of specific lengths in graphs with given degree conditions. A notable related result is the theorem by Thomassen (1983), stating that every graph with minimum degree at least 3 contains a cycle of even length. The approach here draws inspiration from probabilistic methods and the cycle space over $\mathbb{F}_2$.

\section{Lemma Decomposition and Proof Strategy}

\begin{lemma}
Let $G = (V, E)$ be a graph with $\delta(G) \geq 3$. The cycle space of $G$ over $\mathbb{F}_2$ is spanned by the chordless cycles.
\end{lemma}
\begin{proof}
Let $\mathcal{C}$ be the cycle space of $G$. By standard algebraic graph theory, any cycle can be expressed as a symmetric difference of basic cycles. If a cycle $C$ has a chord, we can decompose $C$ into two smaller cycles $C_1$ and $C_2$ such that $C = C_1 \oplus C_2$. Since $G$ is finite, this decomposition must terminate at chordless cycles. Thus, the chordless cycles span $\mathcal{C}$.
\end{proof}

\begin{lemma}
If $G$ is a cubic graph, the length of the longest chordless cycle is bounded from below.
\end{lemma}
\begin{proof}
Consider a maximal path $P = (v_0, v_1, \ldots, v_k)$ in $G$. Since $\delta(G) \geq 3$, $v_0$ must have at least two neighbors in $P$ other than $v_1$, say $v_i$ and $v_j$ with $i < j$. The edge $(v_0, v_j)$ forms a cycle $C$ of length $j+1$. Since $P$ is maximal and $G$ is cubic, topological constraints on the degree enforce a minimum length for such cycles to avoid multi-edges or loops.
\end{proof}

\section{Architecture for Autoformalization (Lean 4)}
The formalization strategy in Lean 4 will define the graph structure using Mathlib's \texttt{SimpleGraph}.

\begin{verbatim}
import Mathlib.Combinatorics.SimpleGraph.Basic
import Mathlib.Data.Nat.Basic

variable {V : Type*} [Fintype V] [DecidableEq V]
variable (G : SimpleGraph V)

def min_degree_ge_three (G : SimpleGraph V) : Prop :=
  \forall v : V, 3 \leq G.degree v

def is_power_of_two (n : \mathbb{N}) : Prop :=
  \exists k : \mathbb{N}, n = 2^k
\end{verbatim}

\end{document}
